% Generated from the audited Markdown source; responsible author: Carptopus.
\documentclass[11pt]{article}
\usepackage[a4paper,margin=28mm]{geometry}
\usepackage[T1]{fontenc}
\usepackage[utf8]{inputenc}
\usepackage{lmodern}
\usepackage{microtype}
\usepackage{amsmath,amssymb,amsthm,mathtools}
\usepackage{needspace}
\usepackage{xcolor}
\usepackage[colorlinks=true,linkcolor=blue!55!black,citecolor=blue!55!black,urlcolor=blue!65!black]{hyperref}
\usepackage{fancyhdr}

\newtheorem{theorem}{Theorem}[section]
\newtheorem{lemma}[theorem]{Lemma}
\newtheorem{corollary}[theorem]{Corollary}
\theoremstyle{remark}
\newtheorem{remark}[theorem]{Remark}
\numberwithin{equation}{section}

\pagestyle{fancy}
\fancyhf{}
\fancyhead[L]{Carptopus}
\fancyhead[R]{Ferrers-polytope face lattice}
\fancyfoot[C]{\thepage}
\setlength{\headheight}{14pt}
\setlength{\parindent}{0pt}
\setlength{\parskip}{0.45em}
\setlength{\emergencystretch}{4em}

\title{The face lattice of the irreducible-Ferrers polytope\\
is a product of triangles}
\author{Carptopus\\\href{mailto:carptopus@163.com}{carptopus@163.com}}
\date{24 August 2026}

\hypersetup{
  pdftitle={The face lattice of the irreducible-Ferrers polytope is a product of triangles},
  pdfauthor={Carptopus},
  pdfsubject={A proof of the product-of-triangles conjecture for irreducible-Ferrers polytopes},
  pdfkeywords={Ferrers diagram, polytope, face lattice, product of triangles, rank-metric code, combinatorial equivalence}
}

\begin{document}
\maketitle

\begin{center}
\fcolorbox{black!35}{black!4}{\parbox{0.88\textwidth}{\centering\small
Public Beta, version 0.1. The mathematical claims have passed a
project-internal zero-trust manuscript audit; external review is pending.}}
\end{center}

\begin{abstract}


Beeloo-Sauerbier Couvée and Neri associated, for every integer $d\geq 3$,
an integral polytope $\mathfrak P_d$ with irreducible Ferrers diagrams and
conjectured that $\mathfrak P_d$ is combinatorially equivalent to a Cartesian
product of $d-2$ triangles. Their construction includes a linearly isomorphic
nonnegative affine slice $\mathfrak L_d$ and a complete description of its
vertices: after grouping the coordinates into $d-2$ triples, each vertex has
exactly one positive coordinate in every triple, and all $3^{d-2}$ possible
support patterns occur. We show that this vertex-support description already
determines the whole face lattice. More generally, a polytope obtained by
intersecting an affine subspace with a nonnegative orthant, whose vertices
realize exactly the one-per-block supports, has the face lattice of a product
of simplices. Applying this observation proves the product-of-triangles
conjecture for every $d\geq 3$. It also gives the complete face enumerator,
simplicity, and the number of facets of $\mathfrak P_d$.

\end{abstract}

\section{Introduction}


Ferrers-diagram rank-metric codes connect algebraic coding theory with the
combinatorics of partitions and convex polytopes. In recent work,
Beeloo-Sauerbier Couvée and Neri \cite{couvee2026} introduced a reduction theory for Ferrers
diagrams in the Etzion--Silberstein conjecture. For every distance parameter
$d\geq 3$, they encoded the irreducible diagram pairs as the integer points of
an integral polytope $\mathfrak P_d$.

The same paper computed the face vectors of $\mathfrak P_d$ for
$3\leq d\leq 7$. These vectors agree with those of the Cartesian product of
$d-2$ triangles, leading to the following conjecture \cite[Conjecture 5.28]{couvee2026}:

\begin{quote}\itshape
For every $d\geq 3$, the polytope $\mathfrak P_d$ is combinatorially
equivalent to the Cartesian product of $d-2$ triangles.
\end{quote}


The purpose of this note is to prove that conjecture. The argument uses no
additional enumeration. The key input is the vertex theorem already proved
in \cite[Proposition 5.25]{couvee2026}. After a sequence of linear isomorphisms, the
polytope is presented as a nonnegative affine slice with coordinates grouped
into triples. Its vertices realize all possible choices of one positive
coordinate per triple. In such a presentation, a face is determined by the
coordinates that vanish on its relative interior, and averaging compatible
vertices realizes every allowable vanishing pattern. Those patterns are
exactly the faces of a product of triangles.

We isolate this argument as a general block-support lemma for products of
simplices. The Ferrers-polytope conjecture is its immediate three-coordinate
specialization. This formulation separates the new observation from the
substantial prior work in \cite{couvee2026}: the definition and integrality of
$\mathfrak P_d$, the affine model $\mathfrak L_d$, and the complete vertex
classification all belong to Beeloo-Sauerbier Couvée and Neri. The new result
is the recovery of the full face lattice from that classification.

Throughout, a \emph{polytope} is a nonempty bounded polyhedron. Two polytopes are
\emph{combinatorially equivalent} if their face lattices, including the empty face,
are isomorphic. We write $\Delta_{r-1}$ for an $(r-1)$-simplex.

\section{Faces of nonnegative affine slices}


We first recall directly, in the form needed below, how faces are encoded in
a nonnegative affine-slice presentation.

Let $A\subseteq\mathbb R^N$ be an affine subspace and let

\nopagebreak[4]
\[
P=A\cap\mathbb R_{\geq 0}^{N}
\]

be a polytope. For a nonempty face $F$ of $P$, choose
$p\in\mathrm{relint}(F)$ and define

\nopagebreak[4]
\[
Z(F)=\{j\in[N]:p_j=0\}.
\]

The set $Z(F)$ is independent of the chosen relative-interior point, and

\nopagebreak[4]
\[
F=P\cap\bigcap_{j\in Z(F)}\{x_j=0\}.
\tag{2.1}
\]

For completeness, let $y$ belong to the right-hand side of (2.1). Every
coordinate of $p$ outside $Z(F)$ is positive, so for sufficiently small
$\varepsilon>0$,

\nopagebreak[4]
\[
w=(1+\varepsilon)p-\varepsilon y
\]

still belongs to $A\cap\mathbb R^N_{\geq0}=P$. Since
$p=(w+\varepsilon y)/(1+\varepsilon)$, the face property and
$p\in\mathrm{relint}(F)$ imply $w,y\in F$. Thus (2.1) holds. Conversely,
an intersection with coordinate hyperplanes as in (2.1), when nonempty, is a
face of $P$.

We now group the coordinates. Let $I$ be a finite index set and, for every
$i\in I$, let $E_i$ be a finite set with $|E_i|\geq2$. The ambient coordinates
are indexed by pairs $(i,e)$ with $i\in I$ and $e\in E_i$.

\Needspace{0.32\textheight}
\begin{lemma}[complete block-support lemma]

Let

\nopagebreak[4]
\[
P=A\cap\mathbb R_{\geq0}^{\sum_{i\in I}|E_i|}
\]

be a polytope. Suppose its vertices are precisely

\nopagebreak[4]
\[
v_\sigma,
\qquad
\sigma\in\prod_{i\in I}E_i,
\]

and their supports satisfy

\nopagebreak[4]
\[
(v_\sigma)_{i,e}>0
\quad\Longleftrightarrow\quad
e=\sigma_i.
\tag{2.2}
\]

Then the face lattice of $P$ is isomorphic to the face lattice of

\nopagebreak[4]
\[
\prod_{i\in I}\Delta_{|E_i|-1}.
\]

\end{lemma}

\begin{proof}

Let $F$ be a nonempty face of $P$. For every $i\in I$, its zero set $Z(F)$
cannot contain all coordinates $(i,e)$ with $e\in E_i$: otherwise every
vertex of $F$ would have zero support in the entire $i$-th block, contradicting
(2.2). Hence

\nopagebreak[4]
\[
S_i(F)=\{e\in E_i:(i,e)\notin Z(F)\}
\]

is nonempty. Equation (2.1) shows that $F$ is uniquely determined by the tuple
$(S_i(F))_{i\in I}$.

It remains to show that every tuple of nonempty subsets occurs. Fix nonempty
$S_i\subseteq E_i$ and set

\nopagebreak[4]
\[
\Sigma_S=\{\sigma\in\prod_{i\in I}E_i:
\sigma_i\in S_i\text{ for every }i\}.
\]

This set is nonempty. Let

\nopagebreak[4]
\[
q_S=\frac{1}{|\Sigma_S|}\sum_{\sigma\in\Sigma_S}v_\sigma.
\]

By (2.2), the coordinate $(i,e)$ of $q_S$ is positive exactly when
$e\in S_i$. Therefore the zero set of $q_S$ is precisely

\nopagebreak[4]
\[
Z_S=\{(i,e):e\notin S_i\}.
\]

The coordinate face

\nopagebreak[4]
\[
F_S=P\cap\bigcap_{(i,e)\in Z_S}\{x_{i,e}=0\}
\]

is nonempty, contains $q_S$, and has no coordinate inequality beyond those in
$Z_S$ active at $q_S$. Thus $q_S\in\mathrm{relint}(F_S)$ and
$Z(F_S)=Z_S$. In particular, distinct tuples $S$ give distinct faces.

We have obtained a bijection between the nonempty faces of $P$ and tuples of
nonempty subsets $S_i\subseteq E_i$. Under this bijection, face containment is
coordinatewise inclusion of the subsets. This is exactly the nonempty face
poset of $\prod_i\Delta_{|E_i|-1}$. Matching the unique empty faces completes
the face-lattice isomorphism.
\end{proof}


\Needspace{0.32\textheight}
\begin{remark}

The hypotheses require the complete and exact support description (2.2). A
vertex count alone does not determine the face lattice. Strict positivity of
the selected coordinate is also essential: without it, averaging compatible
vertices need not realize the intended zero pattern.

The lemma proves only combinatorial equivalence. It does not produce an affine
or a unimodular equivalence between $P$ and the product of simplices.

\end{remark}

\section{The Ferrers polytopes}


We now apply Lemma 2.1 to the polytopes of \cite{couvee2026}. Fix $d\geq3$ and put

\nopagebreak[4]
\[
k=d-2.
\]

Beeloo-Sauerbier Couvée and Neri construct four linearly isomorphic polytopes

\nopagebreak[4]
\[
\mathfrak P_d,
\qquad
\mathfrak P'_d,
\qquad
\mathfrak K_d,
\qquad
\mathfrak L_d.
\tag{3.1}
\]

The first transformation is unimodular \cite[Lemma 5.23]{couvee2026}. The latter two arise
from injective linear maps whose restrictions are linear isomorphisms onto
their images \cite[Lemma 5.24 and the definitions following it]{couvee2026}. Consequently,
all four polytopes have isomorphic face lattices.

The model relevant here is $\mathfrak L_d\subseteq\mathbb R^{3k}$. Write its
coordinates as

\nopagebreak[4]
\[
(x_1,\ldots,x_k,y_1,\ldots,y_k,z_1,\ldots,z_k).
\]

Equation (11) of \cite{couvee2026} presents $\mathfrak L_d$ as the intersection of the
nonnegative orthant with $k$ affine equations. The precise coefficient matrix
of those equations will not be needed here; what matters is that the only
inequalities in this presentation are

\nopagebreak[4]
\[
x_i\geq0,
\qquad
y_i\geq0,
\qquad
z_i\geq0
\qquad(i\in[k]).
\tag{3.2}
\]

Proposition 5.25 of \cite{couvee2026} gives the complete vertex set. For every word

\nopagebreak[4]
\[
\sigma\in\{X,Y,Z\}^{k},
\]

there is exactly one vertex $v_\sigma$, and in the $i$-th coordinate triple

\nopagebreak[4]
\[
(x_i,y_i,z_i)
\]

the coordinate selected by $\sigma_i$ is strictly positive while the other
two coordinates are zero. Conversely, every vertex is obtained in this way.
The strict positivity follows in \cite{couvee2026} from the explicit consecutive-block
formula in the proof of Proposition 5.25.

We can therefore take $I=[k]$ and $E_i=\{X,Y,Z\}$ in Lemma 2.1.

\Needspace{0.32\textheight}
\begin{theorem}

For every integer $d\geq3$,

\nopagebreak[4]
\[
\mathfrak P_d\simeq_{\mathrm{comb}}(\Delta_2)^{d-2}.
\tag{3.3}
\]

In particular, Conjecture 5.28 of \cite{couvee2026} holds for all $d\geq3$.

\end{theorem}

\begin{proof}

The affine-slice presentation (3.2) and the vertex-support theorem just
described satisfy every hypothesis of Lemma 2.1. Hence

\nopagebreak[4]
\[
\mathfrak L_d\simeq_{\mathrm{comb}}(\Delta_2)^k.
\]

The linear isomorphisms in (3.1) preserve face lattices. Substituting
$k=d-2$ proves (3.3).
\end{proof}


\Needspace{0.32\textheight}
\begin{corollary}

Let $k=d-2$. The polytope $\mathfrak P_d$ is a simple polytope of dimension
$2k$, with $3k$ facets and $3^k$ vertices. Every vertex has degree $2k$.

\end{corollary}

\begin{proof}

These properties hold for a product of $k$ two-dimensional simplices and are
preserved by combinatorial equivalence. The vertex count also agrees with
\cite[Theorem 5.20]{couvee2026}.
\end{proof}


\Needspace{0.32\textheight}
\begin{corollary}

If $f_j(\mathfrak P_d)$ denotes the number of $j$-dimensional faces, then

\nopagebreak[4]
\[
\sum_{j=0}^{2k}f_j(\mathfrak P_d)t^j
=(3+3t+t^2)^k.
\tag{3.4}
\]

Equivalently,

\nopagebreak[4]
\[
f_j(\mathfrak P_d)
=\sum_{b=0}^{\lfloor j/2\rfloor}
\frac{k!}{(j-2b)!\,b!\,(k-j+b)!}
3^{k-b},
\tag{3.5}
\]

where a summand is understood to vanish when one of the factorial arguments
is negative.

\end{corollary}

\begin{proof}

The face polynomial of a triangle, with the empty face omitted, is
$3+3t+t^2$. Face dimensions add under Cartesian products, giving (3.4).
To obtain (3.5), choose $b$ factors contributing a two-dimensional face,
$j-2b$ factors contributing an edge, and all remaining factors contributing
a vertex. The edge and vertex choices contribute
$3^{j-2b}3^{k-j+b}=3^{k-b}$.
\end{proof}


\Needspace{0.32\textheight}
\begin{remark}

Formula (3.4) explains, for every $d$, the face vectors computed in \cite{couvee2026} for
$3\leq d\leq7$. These enumerative statements are consequences of the single
face-lattice theorem, not separate independent results.

\end{remark}

\section{Scope and relation to prior work}


Theorem 3.1 is short because the difficult structural input was already
available. The construction of $\mathfrak P_d$, its interpretation in terms
of irreducible Ferrers diagrams, its integrality, and the exact
$3^{d-2}$-vertex classification are all results of \cite{couvee2026}. The present note adds
the observation that their support classification, together with the
nonnegative affine-slice presentation, determines every face rather than only
the vertices.

The result does not settle the Etzion--Silberstein conjecture itself.
The polytopes $\mathfrak P_d$ parametrize the irreducible cases identified in
\cite{couvee2026}, but knowing their face lattices does not construct maximum Ferrers
diagram codes at their integer points. Nor does Theorem 3.1 assert that
$\mathfrak P_d$ is affinely or unimodularly equivalent to a product of
triangles. Such stronger equivalences would require metric or lattice data not
encoded by the face lattice.

As of 24 August 2026, the arXiv record of \cite{couvee2026} contains only version v1 and the
authors' public repository contains only its initial commit. Exact searches
for Conjecture 5.28 and its product-of-triangles formulation found no direct
public proof. This search boundary supports the relevance of the result but
is not a claim of exhaustive global priority.

\section*{Acknowledgement of AI assistance}


The author used OpenAI Codex as an AI-assisted tool for literature discovery,
proof stress testing, computationally checking finite consequences, and
language editing. The author reviewed the mathematical argument and assumes
responsibility for the claims and presentation.

\begin{thebibliography}{9}
\footnotesize
\raggedright

\bibitem{couvee2026}
H.~Beeloo-Sauerbier Couv\'ee and A.~Neri,
\emph{Irreducible Ferrers diagrams in the Etzion--Silberstein conjecture},
arXiv:2604.27868v1 (2026).
\href{https://arxiv.org/abs/2604.27868}{arXiv:2604.27868}.

\end{thebibliography}

\end{document}
